\capitulofrontal{Como usar este ebook}{como-usar}

Este é um ebook \textbf{interativo}: quase tudo o que você vê com cor pode ser clicado.
Ele foi pensado para ser estudado no computador, no tablet ou no celular, mas também
funciona impresso. Antes de começar, vale entender a lógica da coleção.

\secaofrontal{A trilha de cada capítulo}

Todos os capítulos seguem a mesma trilha, que acompanha a forma como o ENEM avalia você:

\begin{center}
\begin{tikzpicture}[
  etapa/.style={rounded corners=5pt,fill=ebookSoft,draw=ebookLight,line width=.7pt,
    minimum height=17mm,text width=25mm,align=center,font=\sffamily\footnotesize},
  seta/.style={-{Stealth[length=2.6mm]},ebookPrim,line width=1.1pt}]
  \node[etapa] (a) at (0,0) {{\color{ebookPrim}\Large\faBookOpen}\\[1mm]\textbf{Teoria}\\exemplos resolvidos};
  \node[etapa] (b) at (3.05,0) {{\color{dicaC}\Large\faLightbulb}\\[1mm]\textbf{Dicas e macetes}\\pegadinhas e resumo};
  \node[etapa] (c) at (6.1,0) {{\color{ebookDark}\Large\faUniversity}\\[1mm]\textbf{Questões do ENEM}\\em ordem de TRI};
  \node[etapa] (d) at (9.15,0) {{\color{ebookAcc}\Large\faFeather*}\\[1mm]\textbf{Inéditas}\\fácil $\to$ muito difícil};
  \node[etapa] (e) at (12.2,0) {{\color{nivelA}\Large\faCheckCircle}\\[1mm]\textbf{Resoluções}\\comentadas};
  \foreach \u/\v in {a/b,b/c,c/d,d/e} \draw[seta] (\u) -- (\v);
\end{tikzpicture}
\end{center}

\begin{enumerate}[leftmargin=1.6em,itemsep=2pt]
  \item \textbf{Teoria.} Explicação do conteúdo com exemplos resolvidos, do simples ao típico de prova.
  \item \textbf{Dicas e macetes.} Estratégias de leitura, atalhos de cálculo, as pegadinhas mais comuns
        e um resumo para revisar na véspera.
  \item \textbf{Questões de edições anteriores.} Questões \emph{oficiais} do ENEM (\anosbanco), reproduzidas
        do caderno original do INEP e organizadas da \textbf{mais fácil para a mais difícil} segundo a
        dificuldade \emph{real} medida pela TRI.
  \item \textbf{Questões inéditas.} Questões originais no estilo do ENEM, também em ordem crescente de
        dificuldade, para você treinar sem “já ter visto a resposta”.
  \item \textbf{Resoluções comentadas.} No fim de cada capítulo, com o passo a passo e a explicação de
        \emph{por que cada alternativa errada está errada} — é ali que você aprende com os distratores.
\end{enumerate}

\secaofrontal{O que significa cada caixa}

\begin{multicols}{2}
\begin{definicao}[exemplo]
Definições e conceitos que você precisa dominar.
\end{definicao}
\begin{dica}
Estratégias de prova e de leitura de enunciados.
\end{dica}
\begin{macete}
Atalhos de cálculo e de raciocínio que economizam tempo.
\end{macete}
\columnbreak
\begin{atencao}
Erros comuns e armadilhas que viram alternativas erradas.
\end{atencao}
\begin{contexto}
Aplicações do conteúdo no mundo real.
\end{contexto}
\begin{formula}[Fórmula-chave]
$\text{o que você precisa saber de cor}$
\end{formula}
\end{multicols}

\secaofrontal{Os selos de dificuldade}

Cada questão traz um selo com quatro barras. Nas questões oficiais, o nível vem do
\textbf{parâmetro $b$ da TRI}, calculado pelo INEP e publicado nos microdados do ENEM; nas
inéditas, é uma estimativa calibrada pelas questões oficiais do mesmo assunto.

\begin{center}
\renewcommand\arraystretch{1.5}
\begin{tabular}{@{}l l l@{}}
\toprule
\textbf{Selo} & \textbf{Parâmetro $b$} & \textbf{Posição aproximada na escala do ENEM} \\
\midrule
\selonivel{1} & $b < 1$ & abaixo de 600 pontos \\
\selonivel{2} & $1 \le b < 2$ & de 600 a 700 pontos \\
\selonivel{3} & $2 \le b < 3$ & de 700 a 800 pontos \\
\selonivel{4} & $b \ge 3$ & acima de 800 pontos \\
\bottomrule
\end{tabular}
\end{center}

O cabeçalho de cada questão oficial mostra ainda o ano, o número da questão no caderno azul, a
\textbf{habilidade} da Matriz de Referência (H1 a H30) e os três parâmetros da TRI ($a$, $b$ e $c$),
explicados no próximo capítulo.

\secaofrontal{Navegação interativa}

\begin{itemize}[leftmargin=1.4em,itemsep=3pt]
  \item[\textcolor{ebookPrim}{\faListUl}] \textbf{Sumário e marcadores.} Clique em qualquer linha do
        \hyperlink{sumario}{sumário}. O painel de marcadores (\emph{bookmarks}) do seu leitor de PDF
        mostra capítulos, seções e cada questão.
  \item[\textcolor{ebookPrim}{\faBookmark}] \textbf{Abas laterais.} As abas numeradas na borda direita
        das páginas levam direto ao início de cada capítulo; a aba destacada é o capítulo atual.
  \item[\textcolor{ebookPrim}{\faHome}] \textbf{Botões do rodapé.} Sumário, início do capítulo,
        resoluções do capítulo e \hyperlink{gabarito-geral}{gabarito geral} — em todas as páginas.
  \item[\textcolor{ebookPrim}{\faExchange*}] \textbf{Questão $\leftrightarrow$ resolução.} O botão
        \chip{ebookPrim}{Resolução comentada \faArrowRight} abre a resolução; o botão
        \chipborda[ebookPrim]{ebookPrim}{\faUndo\ Voltar à questão} traz você de volta.
  \item[\textcolor{ebookPrim}{\faEdit}] \textbf{Cartão-resposta e progresso.} Marque a alternativa
        escolhida em cada questão e registre seu avanço na caixa \emph{Meu progresso} da abertura
        de cada capítulo. Esses campos funcionam no Adobe Acrobat Reader, nos navegadores
        (Chrome, Edge, Firefox), no Foxit e no Okular; salve o arquivo para guardar suas marcações.
        Alguns leitores de celular mostram os campos mas não os salvam.
\end{itemize}

\secaofrontal{Como estudar com este material}

\begin{dica}[A trilha TRI de estudo]
\begin{enumerate}[leftmargin=1.4em,itemsep=1pt,topsep=0pt]
  \item Leia a teoria e refaça os exemplos sem olhar a solução.
  \item Resolva as questões oficiais \textbf{na ordem}: as primeiras consolidam o básico, que é o que
        mais pesa na TRI; as últimas mostram o teto do assunto.
  \item Resolva as inéditas cronometrando: a meta é em torno de 3 minutos por questão.
  \item Confira no gabarito rápido e leia \emph{todas} as resoluções, inclusive as das questões que
        você acertou — a análise dos distratores mostra os erros que a prova espera que você cometa.
  \item Uma semana depois, refaça as que errou e marque \emph{Refiz as que errei}.
\end{enumerate}
\end{dica}
